% Chapter 8: Fermat's Last Theorem and the Modularity Theorem
\let\LocalMacro\relax

\chapter{Fermat's Last Theorem and the Modularity Theorem}

\section{The Problem}

\subsection{Informal}

\textbf{Background.} A whole number raised to a power greater than 2 can sometimes be split into two other whole numbers raised to that same power. For example, $3^2 + 4^2 = 5^2$ works for squares. Fermat's Last Theorem asks whether this is possible for higher powers.

\textbf{Given.} You are given any whole number $n$ greater than 2, and you are allowed to search for nonzero whole numbers $x$, $y$, and $z$.

\textbf{Goal.} Determine whether the equation $x^n + y^n = z^n$ has any solutions in nonzero whole numbers, for every exponent $n > 2$.

\textbf{Constraints.} All of $x$, $y$, and $z$ must be positive whole numbers. The exponent $n$ must be greater than 2. The result must hold for every such $n$, with no exceptions.

\textbf{Example.} For $n = 2$: $3^2 + 4^2 = 5^2$ works (Pythagorean triples). For $n = 3$: is there a way to write two perfect cubes that add to a third perfect cube? Fermat claimed there is not, and that the same is true for every higher power. The theorem says the only way to split a perfect power into two perfect powers of the same degree is with squares.

\subsection{Formal}
The equation
\begin{equation}
x^n + y^n = z^n
\end{equation}
has no nonzero integer solutions for $n > 2$.

More deeply, the Taniyama--Shimura--Weil conjecture (now the Modularity Theorem) asserts:

\begin{conjecture}[Modularity Conjecture]
Every elliptic curve over $\mathbb{Q}$ is modular, meaning its $L$-function arises from a modular form.
\end{conjecture}

If this conjecture holds, then Fermat's Last Theorem follows, because a hypothetical counterexample $a^n + b^n = c^n$ would yield an elliptic curve (the Frey curve $y^2 = x(x - a^n)(x + b^n)$) that could not be modular, contradicting the conjecture.

\subsection{Historical Context}
Fermat's note in the margin of Diophantus's \emph{Arithmetica} around 1637 became the most famous claim in the history of mathematics. For over 350 years, the theorem resisted all attempts at proof, attracting efforts from Euler, Gauss, Kummer, and countless others. The eventual resolution required an entirely new branch of mathematics connecting elliptic curves and modular forms.

\subsection{What Was Known}
For over 350 years, mathematicians chipped away at special cases. By the 1990s, the theorem had been verified by computer for all exponents up to about four million. The critical breakthrough occurred in the 1980s when Gerhard Frey and Ken Ribet proved a surprising link: if Fermat's Last Theorem were false, it would imply the existence of an elliptic curve so strange that it could not be modular. This meant that proving the Taniyama--Shimura--Weil conjecture (the Modularity Theorem) would automatically prove Fermat's Last Theorem.

\subsection{Why It Matters}
Fermat's Last Theorem is not just a curiosity about whole numbers. Proving it required building an entirely new bridge between two distant areas of mathematics---elliptic curves and modular forms---that has since become one of the most powerful tools in number theory. The techniques developed along the way have been applied to problems in cryptography, coding theory, and beyond.

\subsection{Simplified Version}
The theorem is easiest to understand through examples: Fermat himself proved the case for $n=4$ using a technique called infinite descent. Leonhard Euler settled $n=3$ in 1770, and others handled $n=5$ shortly after. These special cases suggested the theorem was true, but proving it for \emph{all} possible exponents at once required a mathematical leap that would take centuries to achieve.

\section{Mathematical Theory}


\subsection{Prerequisites}
This chapter assumes familiarity with:
\begin{itemize}
\item Elementary number theory (congruences, factorization)
\item Basic algebraic geometry (elliptic curves)
\item Complex analysis (modular forms basics)
\end{itemize}

\subsection{Elliptic Curves}

An elliptic curve is a smooth projective curve of genus 1 with a rational point. In Weierstrass form:

\begin{equation}
y^2 = x^3 + ax + b
\end{equation}

where $4a^3 + 27b^2 \neq 0$. Elliptic curves form an abelian group under a geometric addition law.

\subsection{Modular Forms}

A modular form is a holomorphic function on the upper half-plane satisfying a transformation law under $SL_2(\mathbb{Z})$:

\begin{equation}
f\left(\frac{az+b}{cz+d}\right) = (cz+d)^k f(z)
\end{equation}

for integers $a,b,c,d$ with $ad-bc = 1$. The Fourier expansion of a modular form encodes deep arithmetic information.

\subsection{The Taniyama--Shimura--Weil Conjecture}

\begin{conjecture}[Modularity Conjecture]
Every elliptic curve over $\mathbb{Q}$ is modular, meaning its L-function arises from a modular form.
\end{conjecture}

This conjecture connected two seemingly unrelated areas: elliptic curves (geometry) and modular forms (analysis).

\subsection{The Frey Curve and Ribet's Theorem}

In 1984, Gerhard Frey observed that if Fermat's Last Theorem were false, one could construct an elliptic curve (now called the Frey curve) with bizarre properties \cite{frey1984}:

\begin{equation}
y^2 = x(x - a^n)(x + b^n)
\end{equation}

where $a^n + b^n = c^n$. This curve would be so unusual that it could not possibly be modular.

\begin{theorem}[Ribet, 1986 \cite{ribet1990}]
If the Modularity Conjecture is true, then Fermat's Last Theorem is true.
\end{theorem}

This was the crucial link: proving modularity for all elliptic curves would settle Fermat's Last Theorem.

\section{The Solution}

\subsection{The Big Idea}
Frey and Ribet showed that a counterexample to Fermat's Last Theorem would produce the Frey curve $y^2 = x(x-a^n)(x+b^n)$, an elliptic curve so pathological it could not be modular. So proving the Modularity Conjecture---every elliptic curve over $\mathbb{Q}$ is modular---would settle FLT. Wiles proved this for semistable elliptic curves using two innovations. First, \textbf{deformation theory}: instead of studying a single elliptic curve, study the \emph{family} of all Galois representations that ``look like'' it. Concretely, given an elliptic curve $E/\mathbb{Q}$ and a prime $\lambda$, there is a Galois representation $\rho_{E,\lambda}: \text{Gal}(\overline{\mathbb{Q}}/\mathbb{Q}) \to \text{GL}_2(\mathbb{Z}_\lambda)$. The \textbf{deformation ring} $R$ parameterizes all Galois representations that are ``deformations'' of $\rho_{E,\lambda}$---representations that reduce to $\rho_{E,\lambda}$ mod $\lambda$ but can vary freely over $\mathbb{Z}_\lambda$. The \textbf{Hecke algebra} $T$ is the subring of $\text{End}(J_0(N))$ generated by Hecke operators---it parameterizes the Galois representations that come from modular forms. The modularity question becomes: is $R = T$? If so, every deformation (every representation that ``looks like'' $\rho_{E,\lambda}$) comes from a modular form. Second, the \textbf{Taylor--Wiles patching argument}: to prove $R \to T$ is surjective, add auxiliary primes to kill obstructions, pass to an inverse limit (the ``patched module'' $M_\infty$) that becomes free over a power series ring, and read off surjectivity. The patched module is constructed so that its structure is transparent---freeness over the power series ring forces $R \to T$ to be surjective, hence every Galois representation in the deformation family is modular. Breuil--Conrad--Diamond--Taylor (2001) extended this to all elliptic curves.

\subsection{What Was Stopping Progress}
Before Wiles, modularity was known only for special elliptic curves---those with unusual symmetry (complex multiplication) or other rare algebraic properties. Ordinary elliptic curves, the vast majority, were completely out of reach. It was like having a machine that could identify a few exotic species of fish but had no theory covering the ocean as a whole. The general problem---proving that *every* elliptic curve over the rationals is modular---required a tool that worked uniformly, without assuming any special structure. No such tool existed.

\subsection{The Breakthrough}
Wiles's key idea was *deformation theory*: instead of studying a single elliptic curve in isolation, study the *family* of all Galois representations that ``look like'' it, and prove that this entire family comes from modular forms. He turned modularity into a question about whether one algebraic ring maps onto another. The Taylor--Wiles patching argument then showed this map is surjective by a clever bookkeeping trick: add auxiliary data to kill off obstructions, pass to an inverse limit, and watch the obstruction vanish. It was like proving a bridge is strong enough by showing every possible load it could face is supported.

\subsection{How It Works}

Before Wiles, the Modularity Conjecture---the statement that every elliptic curve over $\mathbb{Q}$ is modular---was wide open in full generality. Frey (1984) and Ribet (1986) had shown that Fermat's Last Theorem follows from modularity: a counterexample to FLT would produce the Frey curve, an elliptic curve so ``pathological'' that it could not be modular, contradicting the conjecture. But this merely \emph{reduced} FLT to the Modularity Conjecture. The conjecture itself---connecting Galois representations attached to arbitrary elliptic curves with modular forms---had no proof for even the semistable case, let alone all elliptic curves.

Before Wiles, modularity was known only for special classes of elliptic curves: those with complex multiplication (CM), or curves whose Galois representations happened to factor through known Hecke characters. These methods exploited special algebraic structure that generic elliptic curves do not possess. The general problem required a tool that could handle \emph{all} elliptic curves uniformly---proving that a Galois representation ``comes from'' a modular form, without assuming any special properties of the curve. Classical approaches (explicit construction of modular forms, direct comparison of $L$-functions) did not scale to this level of generality.

Wiles's breakthrough rested on a fundamentally new strategy:
\begin{enumerate}
\item \textbf{Deformation theory of Galois representations:} Rather than studying a single elliptic curve in isolation, Wiles considered the \emph{family} of all Galois representations that ``look like'' the one attached to a given elliptic curve. He formulated modularity as a question about whether a certain deformation ring maps surjectively onto a Hecke algebra.
\item \textbf{The Taylor--Wiles patching argument:} To prove this surjectivity, Wiles and Taylor introduced a clever patching construction: they chose auxiliary primes to kill unwanted Selmer groups, then passed to an inverse limit (the ``patched module'') that became free over a power series ring, from which the surjectivity followed. This patching method has since become a standard technique in number theory.
\end{enumerate}

Imagine you want to prove that every lock in a city can be opened by a key from a single key-ring. Frey and Ribet showed that if there were a lock with no matching key, it would produce an ``impossible'' lock---so proving every lock is on the ring suffices. Previous mathematicians had shown this for special locks (those with unusual ward patterns), but ordinary locks remained. Wiles's idea was revolutionary: instead of checking each lock one by one, he studied the \emph{space of all possible locks} and showed that the key-ring ``reaches'' all of them. He did this by studying how locks deform when you make small changes (deformation theory) and then showing, via a clever bookkeeping trick (patching), that the key-ring covers every deformation. Once every lock is accounted for, Fermat's Last Theorem falls as a corollary.

\bigskip

\subsection{Wiles's Proof (1993--1995)}

Andrew Wiles, a British mathematician at Princeton, announced a proof of the Modularity Conjecture for semistable elliptic curves in 1993 \cite{wiles1995}. This immediately implied Fermat's Last Theorem.

The proof used:

\begin{enumerate}
\item \textbf{Galois representations}: Associating Galois representations to elliptic curves and modular forms.
\item \textbf{Deformation theory}: Studying how Galois representations vary in families.
\item \textbf{Iwasawa theory}: Deep results about the arithmetic of elliptic curves over $\mathbb{Z}_p$-extensions.
\item \textbf{Kolyvagin--Flach Euler systems}: New tools for bounding Selmer groups.
\end{enumerate}

\begin{highlightbox}
Wiles worked in secret for seven years on this proof. He told no one except his wife. On June 21, 1993, he gave a three-hour lecture at the Isaac Newton Institute in Cambridge, England, revealing the proof. The audience sat in stunned silence when he finished.
\end{highlightbox}

\subsection{The Flaw and the Fix}

After peer review, a gap was found in the proof. Wiles and his former student Richard Taylor spent nearly a year trying to fix it. In September 1994, they realized that combining two different approaches to the deformation theory would work. The corrected proof was published in 1995 in two papers:

\begin{itemize}
\item Wiles (1995): ``Modular elliptic curves and Fermat's Last Theorem'' \cite{wiles1995}
\item Taylor \& Wiles (1995): ``Ring-theoretic properties of certain Hecke algebras'' \cite{taylorwiles1995}
\end{itemize}

\subsection{The Complete Proof}

Breuil, Conrad, Diamond, and Taylor (2001) extended Wiles's methods to prove the full Modularity Conjecture for all elliptic curves over $\mathbb{Q}$ \cite{bcdt2001}, completing the proof of Fermat's Last Theorem.

\subsection{After the Proof}

The Fermat--Euler theorem and the full Modularity Theorem are both completely settled. Wiles proved modularity for semistable elliptic curves, implying Fermat's Last Theorem, and Breuil--Conrad--Diamond--Taylor (2001) extended this to all elliptic curves over $\mathbb{Q}$. There is nothing left to prove: Fermat's Last Theorem is a theorem, and the Taniyama--Shimura--Weil conjecture is a theorem. The proof is unconditional and has been thoroughly verified.

The techniques that Wiles developed---deformation theory of Galois representations, the Taylor--Wiles patching argument, and the analysis of Hecke algebras---have become central tools in modern number theory. They are the foundation for major advances in the Langlands program, which predicts deep connections between automorphic forms, Galois representations, and algebraic geometry. The patching method in particular has been generalized and applied to an extraordinary range of problems: from level-lowering results to the study of Selmer groups, from the Fontaine--Mazur conjecture to the construction of Galois representations with prescribed local behavior.

Several major questions remain open. The Langlands program itself is far from complete: while the modularity theorem handles elliptic curves over $\mathbb{Q}$, the corresponding questions for abelian varieties of higher dimension, for curves of higher genus, and for number fields other than $\mathbb{Q}$ are largely open. The modularity of higher-dimensional abelian varieties---are they all ``automorphic'' in some sense?---is one of the central open problems in the field. Beyond the Langlands program, the deformation-theoretic methods Wiles introduced have not yet led to a proof of major open conjectures like the Birch and Swinnerton-Dyer conjecture, though they provide essential partial results.

\section{Impact}

\begin{itemize}
\item \textbf{Number theory}: Wiles's methods revolutionized the study of Galois representations and deformation theory.
\item \textbf{Langlands program}: The proof confirmed deep connections between different areas of mathematics predicted by the Langlands program.
\item \textbf{Mathematical culture}: Wiles's story captured the public imagination and demonstrated the power of sustained dedication to a single problem.
\end{itemize}

\section{Exercises}

\begin{exercise}[Basic]
Prove that $x^4 + y^4 = z^4$ has no nonzero integer solutions using Fermat's method of infinite descent. \emph{Hint: If $x^4 + y^4 = z^4$, then $(x^2)^2 + (y^2)^2 = (z^2)^2$, giving a Pythagorean triple.}
\end{exercise}

\begin{exercise}[Intermediate]
Show that the elliptic curve $y^2 = x^3 + ax + b$ is nonsingular if and only if $4a^3 + 27b^2 \neq 0$. \emph{Hint: Compute the partial derivatives and find where they simultaneously vanish.}
\end{exercise}

\begin{exercise}[Challenge]
Explain why Ribet's theorem implies that proving modularity for semistable elliptic curves implies Fermat's Last Theorem. \emph{Hint: A Frey curve from a Fermat counterexample would be semistable.}
\end{exercise}


\begin{exercise}[Computational]
Write a Python/SageMath script to explore a concept from this chapter. See the appendix for starter code.
\end{exercise}

\section{Timeline}

\begin{tabularx}{\linewidth}{lX}
\toprule
\textbf{Year} & \textbf{Event} \\ \midrule
1637 & Fermat writes his marginal note \\ 
1753 & Euler proves $n = 3$ \cite{euler1783} \\ 
1825 & Dirichlet and Legendre prove $n = 5$ \\ 
1839 & Lam\'e proves $n = 7$ \\ 
1984 & Frey observes the connection to elliptic curves \cite{frey1984} \\ 
1986 & Ribet proves the connection to modularity \cite{ribet1990} \\ 
1993 & Wiles announces the proof \cite{wiles1995} \\ 
1994 & Wiles and Taylor fix the gap \cite{taylorwiles1995} \\ 
1995 & The complete proof is published \cite{wiles1995,taylorwiles1995} \\ 
2001 & Breuil--Conrad--Diamond--Taylor prove full modularity \cite{bcdt2001} \\ \bottomrule
\end{tabularx}

\section{Citations}

\begin{enumerate}
\item Wiles, A. (1995). ``Modular elliptic curves and Fermat's Last Theorem.'' Annals of Mathematics, 141(3), 443--551.
\item Taylor, R., \& Wiles, A. (1995). ``Ring-theoretic properties of certain Hecke algebras.'' Annals of Mathematics, 141(3), 553--572.
\item Breuil, C., Conrad, B., Diamond, F., \& Taylor, R. (2001). ``On the modularity of elliptic curves over $\mathbb{Q}$.'' Inventiones Mathematicae, 145(3), 493--541.
\item Ribet, K. (1990). ``On modular representations of $Gal(\overline{\mathbb{Q}}/\mathbb{Q})$ arising from modular forms.'' Inventiones Mathematicae, 100(1), 43--103.
\item Frey, G. (1984). ``Links between stable elliptic curves and certain Diophantine equations.'' Annales Universitatis Saraviensis.
\item Euler, L. (1783). ``De resolutione formulae arithmeticae.'' Novi Commentarii Academiae Scientiarum Petropolitanae.
\end{enumerate}